\documentclass[12pt]{book}
\usepackage{amsmath}
\usepackage{amssymb}
\usepackage{geometry}
\geometry{margin=1in}

\title{Volume 12: Computation, Code, and Simulation Systems \\ Book 02: Simulations and Visualizations}
\author{James C. Harvey}
\date{December 2025}

\begin{document}
\maketitle
\tableofcontents

\chapter{Simulation as Geometric Experiment}

\section*{Abstract}
Simulations in SDT are geometric experiments. They propagate flux and boundary constraints through
numerical forms of the master equation and state operators.

\section{Solar System and Galactic Models}
Orbital and galactic simulations test the pressure-gradient framework at large scales and validate
occlusion-driven rotation curves.

\section{Visualization Systems}
Visualization is not aesthetic garnish; it is an explicit rendering of flux topology and boundary
geometry. Rendered images are therefore computational theorems.

\chapter{Numerical Methods}

\section*{Abstract}
Simulations use numerical integration of pressure gradients and occlusion operators. This chapter
describes the method classes.

\section{Integrator Classes}
Matrix flow is advanced with explicit and semi-implicit schemes that preserve gradient structure:
\begin{equation}
\mathbf{v}_{t+\Delta t} = \mathbf{v}_t + \Delta t\, \mathcal{F}(\nabla \Pi, \nu).
\end{equation}

\chapter{Simulation Pipelines}

\section*{Abstract}
Simulations follow a pipeline from geometry input to rendered output. Each stage is documented for
reproducibility.

\section{Pipeline Stages}
Input geometry $\rightarrow$ state vector construction $\rightarrow$ operator evaluation $\rightarrow$
integration $\rightarrow$ visualization.

\chapter{Validation Outputs}

\section*{Abstract}
Simulation outputs are validated against benchmark artifacts. This chapter explains the comparison
protocol.

\section{Output Comparison}
Numeric outputs are compared to benchmark tables without adjustment to constants or geometry.
\end{document}
